\documentclass{article}
\usepackage{graphicx} % Required for inserting images
\usepackage{booktabs}
\title{Lesson 2 Challenge - Calculus}
\author{William Wang}
\date{August 2026}

\begin{document}

\maketitle

\section{}
We have \(g(x) = \lfloor x \rfloor\), then 
\[\lim_{x\to\frac{3}{2}} g(x) = 1\]
since no matter which way we approach \(\frac{3}{2}\) from it will always be less than \(2\) and greater than \(1\) meaning the floor function will return the same result. However when we have 
\[\lim_{x\to2}g(x)\]
If we approach \(x \to 2^-\) it will be less than $2$, hence the floor will return 1 but if we approach \(x \to 2^+\), then it will be greater than 2, hence the floor will return 2. Therefore, since there are $2$ values, it is undefined for \(x \to 2\). 
\section{}
We can construct a table of values. 
\begin{table}[ht]
\centering
\begin{tabular}{cc}
\toprule
$x$ & $g \circ f(x)$\\
\midrule
$1.51$    & $9$ \\
$1.501$   & $9$ \\
$1.5001$  & $9$ \\
$1.50001$ & $9$ \\

$1.49$& $8$ \\
$1.499$   & $8$ \\
$1.4999$    & $8$ \\
$1.49999$     & $8$ \\
\bottomrule
\end{tabular}
\end{table}
We can see that depending on which way you approach \(\frac{3}{2}\), the value of \(g\circ f(x)\) changes, hence the limit is undefined. 
\section{}
Using a table of values,
\begin{table}[ht]
\centering
\begin{tabular}{cc}
\toprule
$x$ & $f \circ g(x)$\\
\midrule
$1.51$    & $4$ \\
$1.501$   & $4$ \\
$1.5001$  & $4$ \\
$1.50001$ & $4$ \\

$1.49$& $4$ \\
$1.499$   & $4$ \\
$1.4999$    & $4$ \\
$1.49999$     & $4$ \\
\bottomrule
\end{tabular}
\end{table}
From the table of values we can see that 
\[\lim_{x\to \frac{3}{2}} f\circ g(x) = 4\] Hence the limit exists. 

\section{}
They are not enough because the order in which you do them matters. The output of one of the function is the input of the second, meaning it could be a value where the limit if undefined for that part. 

\section{}
If we have 
\[\lim_{x\to c} f (g(x))\]
then \(\lim_{x\to c} g(x)\) has to be defined and we also need a function \(f(x)\) such that there exists a limit at every point on \(f(x)\). Then we can pull the limit into the inner function such that 
\[\lim_{x\to c} f(g(x)) = f(\lim_{x\to c}g(x))\]

\end{document}
